\documentclass{article}
\usepackage[T1]{fontenc}
\usepackage{lmodern}
\usepackage{graphicx}
\usepackage{amsmath,amssymb}
\usepackage[a4paper,margin=2cm]{geometry}
\usepackage[absolute,overlay]{textpos}
\setlength{\TPHorizModule}{1mm}
\setlength{\TPVertModule}{1mm}
\usepackage[indonesian]{babel}
\usepackage{hyperref}
\setlength{\emergencystretch}{2em}
% unit: Halaman 16

\begin{document}

% === Halaman 16 (python/native) ===

• Alice mendekripsi cipherteks C = 0328, 0301, 2653, 2986, 1164 dengan 

menggunakan kunci privatnya, yaitu  d = 1019, n = 3337:


c1 = 0328 →m1 = 3281019 mod 3337 = 704 = 0704 

c2 = 0301 →m2 = 3011019 mod 3337 = 1111 

c3 = 2653 →m3 = 26531019 mod 3337 = 1400

c4 = 2986 →m4 = 29861019 mod 3337 = 1108

c5 = 1164 →m5 = 11641019 mod 3337 = 204

• Alice memperoleh kembali plainteks dari Bob 

m = 07041111140011080204 

   yang dikodekan kembali (A = 00, B = 01, C = 02, …, Z = 25) menjadi 


M = HELLO ALICE 

Dekripsi: m = cd mod n

16

\end{document}
